% Generated draft: requires mathematical and source review.
\documentclass[11pt]{article}
\usepackage[margin=25mm]{geometry}
\usepackage{fontspec}
\setmainfont{Latin Modern Roman}
\usepackage{amsmath,amssymb,mathtools,unicode-math}
\setmathfont{Latin Modern Math}
\usepackage{enumitem,needspace,xurl,hyperref}
\setlength{\parindent}{0pt}
\setlength{\parskip}{5pt}
\setlength{\emergencystretch}{3em}
\newcommand{\N}{\mathbb{N}}
\newcommand{\Z}{\mathbb{Z}}
\newcommand{\Q}{\mathbb{Q}}
\newcommand{\R}{\mathbb{R}}
\newcommand{\C}{\mathbb{C}}
\newcommand{\F}{\mathbb{F}}
\newcommand{\A}{\mathbb{A}}
\newcommand{\PP}{\mathbb{P}}
\newcommand{\E}{\mathbb{E}}
\newcommand{\eps}{\varepsilon}
\newcommand{\dd}{\,\mathrm{d}}
\begin{document}
\section*{UnsolvedMath ID: 30000659}
\textbf{Source problem code:} catalogue entry
\subsection{Definitions and mathematical statement}
Let \(X\) be a smooth affine complex variety of complex dimension three, whose underlying complex-analytic space \(X(\C)\) is contractible. Let the additive algebraic group \(\mathbb G_a=(\C,+)\) act algebraically and nontrivially, and let
\[
\pi:X\longrightarrow S=X//\mathbb G_a
=\operatorname{Spec}(\mathcal O(X)^{\mathbb G_a})
\]
be the quotient in the source. Finite generation and the precise quotient category are implicit in that notation, not separately given as a theorem here.

Let \(\Gamma\subset S\) be the stated minimal closed curve such that, with \(U=S\setminus\Gamma\), there is an isomorphism over \(U\)
\[
\pi^{-1}(U)\cong U\times\A^1_\C,
\]
commuting with the maps to \(U\). The source does not explicitly demand a \(\mathbb G_a\)-equivariant isomorphism; that stronger trivialization is not imposed silently. Nor is the group action assumed free or the quotient surface smooth.

The question is whether
\[
\Gamma(\C)\text{ connected}\quad\Longrightarrow\quad
\Gamma(\C)\text{ contractible}
\]
for every indicated action and every intended minimal exceptional curve. Contractibility uses the ordinary complex-analytic topology, meaning homotopy equivalence to a point, not \(\A^1\)-homotopy or merely vanishing rational homology.

Minimality must be defined carefully: the usual intended reading is inclusion-minimal among admissible closed exceptional sets for which the displayed product trivialization holds on the complement. The source does not say whether competitors may be empty or zero dimensional, whether only subcurves are allowed, or whether the minimal set is unique. Those choices materially change the hypothesis. The treatment of empty \(\Gamma\), reductions, and connectedness conventions likewise needs confirmation; nilpotent structure is not detected by the underlying topological contractibility question. The existence of such a minimal curve is part of the stated setup, not proved by the draft.

The supplied literature triage checked 2026-08-21 reports no authoritative complete resolution of the extracted claim, without a separate quotient-curve theorem or minimality definition.
\subsection{Short English statement}
For the quotient of a contractible affine threefold by an additive-group action, must a connected minimal exceptional curve be topologically contractible? The product trivialization is over its complement, and the precise minimal-exceptional-set convention is essential.
\subsection{Sources}
\begin{itemize}
\item Local UnsolvedMath registry, permanent ID 30000659, catalogue entry, \emph{Contractibility of Quotient Curves from Additive-Group Actions}.
\item \emph{Affine Algebraic Geometry} (2007), Oberwolfach Reports, pp. 5--82, DOI \url{https://doi.org/10.4171/owr/2007/01}.
\item Supplied literature triage checked 2026-08-21 cites only that report; no separate quotient, trivialization, or minimality theorem is supplied.
\end{itemize}
\end{document}
